\section{Effective Cartier divisors}
\label{section-effective-Cartier-divisors}

\noindent
We define the notion of an effective Cartier divisor before any other type
of divisor.

\begin{definition}
\label{definition-effective-Cartier-divisor}
Let $S$ be a scheme.
\begin{enumerate}
\item A {\it locally principal closed subscheme} of $S$ is a closed subscheme
whose sheaf of ideals is locally generated by a single element.
\item An {\it effective Cartier divisor} on $S$ is a closed subscheme
$D \subset S$ whose ideal sheaf $\mathcal{I}_D \subset \mathcal{O}_S$
is an invertible $\mathcal{O}_S$-module.
\end{enumerate}
\end{definition}

\noindent
Thus an effective Cartier divisor is a locally principal closed subscheme,
but the converse is not always true. Effective Cartier divisors are closed
subschemes of pure codimension $1$ in the strongest possible sense. Namely
they are locally cut out by a single element which is a nonzerodivisor.
In particular they are nowhere dense.

\begin{lemma}
\label{lemma-characterize-effective-Cartier-divisor}
Let $S$ be a scheme.
Let $D \subset S$ be a closed subscheme.
The following are equivalent:
\begin{enumerate}
\item The subscheme $D$ is an effective Cartier divisor on $S$.
\item For every $x \in D$ there exists an affine open neighbourhood
$\Spec(A) = U \subset S$ of $x$ such that
$U \cap D = \Spec(A/(f))$ with $f \in A$ a nonzerodivisor.
\end{enumerate}
\end{lemma}

\begin{proof}
Assume (1).  For every $x \in D$ there exists an affine open neighbourhood
$\Spec(A) = U \subset S$ of $x$ such that
$\mathcal{I}_D|_U \cong \mathcal{O}_U$. In other words, there exists
a section $f \in \Gamma(U, \mathcal{I}_D)$ which freely generates the
restriction $\mathcal{I}_D|_U$. Hence $f \in A$, and the multiplication
map $f : A \to A$ is injective. Also, since $\mathcal{I}_D$ is
quasi-coherent we see that $D \cap U = \Spec(A/(f))$.

\medskip\noindent
Assume (2). Let $x \in D$. By assumption there exists an affine open
neighbourhood $\Spec(A) = U \subset S$ of $x$ such that
$U \cap D = \Spec(A/(f))$ with $f \in A$ a nonzerodivisor.
Then $\mathcal{I}_D|_U \cong \mathcal{O}_U$ since it is equal to
$\widetilde{(f)} \cong \widetilde{A} \cong \mathcal{O}_U$.
Of course $\mathcal{I}_D$ restricted to the open subscheme
$S \setminus D$ is isomorphic to $\mathcal{O}_{S \setminus D}$.
Hence $\mathcal{I}_D$ is an invertible $\mathcal{O}_S$-module.
\end{proof}

\begin{lemma}
\label{lemma-complement-locally-principal-closed-subscheme}
Let $S$ be a scheme. Let $Z \subset S$ be a locally principal closed
subscheme. Let $U = S \setminus Z$. Then $U \to S$ is an affine morphism.
\end{lemma}

\begin{proof}
The question is local on $S$, see
Morphisms, Lemmas \ref{morphisms-lemma-characterize-affine}.
Thus we may assume $S = \Spec(A)$ and $Z = V(f)$ for some $f \in A$.
In this case $U = D(f) = \Spec(A_f)$ is affine hence $U \to S$ is affine.
\end{proof}

\begin{lemma}
\label{lemma-complement-effective-Cartier-divisor}
Let $S$ be a scheme. Let $D \subset S$ be an effective Cartier divisor.
Let $U = S \setminus D$. Then $U \to S$ is an affine morphism and $U$
is scheme theoretically dense in $S$.
\end{lemma}

\begin{proof}
Affineness is Lemma \ref{lemma-complement-locally-principal-closed-subscheme}.
The density question is local on $S$, see
Morphisms, Lemma \ref{morphisms-lemma-characterize-scheme-theoretically-dense}.
Thus we may assume $S = \Spec(A)$ and $D$ corresponding to the
nonzerodivisor $f \in A$, see
Lemma \ref{lemma-characterize-effective-Cartier-divisor}.
Thus $A \subset A_f$ which implies that $U \subset S$ is
scheme theoretically dense, see
Morphisms, Example \ref{morphisms-example-scheme-theoretic-closure}.
\end{proof}

\begin{lemma}
\label{lemma-effective-Cartier-makes-dimension-drop}
Let $S$ be a scheme.
Let $D \subset S$ be an effective Cartier divisor.
Let $s \in D$.
If $\dim_s(S) < \infty$, then $\dim_s(D) < \dim_s(S)$.
\end{lemma}

\begin{proof}
Assume $\dim_s(S) < \infty$.
Let $U = \Spec(A) \subset S$ be an affine open neighbourhood
of $s$ such that $\dim(U) = \dim_s(S)$ and such that $D = V(f)$
for some nonzerodivisor $f \in A$ (see
Lemma \ref{lemma-characterize-effective-Cartier-divisor}).
Recall that $\dim(U)$ is the Krull dimension of the ring $A$
and that $\dim(U \cap D)$ is the Krull dimension of the ring $A/(f)$.
Then $f$ is not contained in any minimal prime of $A$.
Hence any maximal chain of primes in $A/(f)$, viewed as a chain
of primes in $A$, can be extended by adding a minimal prime.
\end{proof}

\begin{definition}
\label{definition-sum-effective-Cartier-divisors}
Let $S$ be a scheme. Given effective Cartier divisors
$D_1$, $D_2$ on $S$ we set $D = D_1 + D_2$ equal to the
closed subscheme of $S$ corresponding to the quasi-coherent
sheaf of ideals
$\mathcal{I}_{D_1}\mathcal{I}_{D_2} \subset \mathcal{O}_S$.
We call this the {\it sum of the effective Cartier divisors
$D_1$ and $D_2$}.
\end{definition}

\noindent
It is clear that we may define the sum $\sum n_iD_i$ given
finitely many effective Cartier divisors $D_i$ on $X$
and nonnegative integers $n_i$.

\begin{lemma}
\label{lemma-sum-effective-Cartier-divisors}
The sum of two effective Cartier divisors is an effective
Cartier divisor.
\end{lemma}

\begin{proof}
Omitted. Locally $f_1, f_2 \in A$ are nonzerodivisors, then also
$f_1f_2 \in A$ is a nonzerodivisor.
\end{proof}

\begin{lemma}
\label{lemma-difference-effective-Cartier-divisors}
Let $X$ be a scheme.
Let $D, D'$ be two effective Cartier divisors on $X$.
If $D \subset D'$ (as closed subschemes of $X$), then
there exists an effective Cartier divisor $D''$ such
that $D' = D + D''$.
\end{lemma}

\begin{proof}
Omitted.
\end{proof}

\begin{lemma}
\label{lemma-sum-closed-subschemes-effective-Cartier}
Let $X$ be a scheme. Let $Z, Y$ be two closed subschemes of $X$
with ideal sheaves $\mathcal{I}$ and $\mathcal{J}$. If $\mathcal{I}\mathcal{J}$
defines an effective Cartier divisor $D \subset X$, then $Z$ and $Y$
are effective Cartier divisors and $D = Z + Y$.
\end{lemma}

\begin{proof}
Applying Lemma \ref{lemma-characterize-effective-Cartier-divisor} we obtain
the following algebra situation: $A$ is a ring, $I, J \subset A$
ideals and $f \in A$ a nonzerodivisor such that $IJ = (f)$.
Thus the result follows from
Algebra, Lemma \ref{algebra-lemma-product-ideals-principal}.
\end{proof}

\begin{lemma}
\label{lemma-sum-effective-Cartier-divisors-union}
Let $X$ be a scheme. Let $D, D' \subset X$ be effective Cartier divisors
such that the scheme theoretic intersection $D \cap D'$ is an effective
Cartier divisor on $D'$. Then $D + D'$ is the scheme theoretic
union of $D$ and $D'$.
\end{lemma}

\begin{proof}
See Morphisms, Definition
\ref{morphisms-definition-scheme-theoretic-intersection-union}
for the definition of scheme theoretic intersection and union.
To prove the lemma working locally
(using Lemma \ref{lemma-characterize-effective-Cartier-divisor})
we obtain the following algebra problem: Given a ring $A$
and nonzerodivisors $f_1, f_2 \in A$ such that $f_1$ maps
to a nonzerodivisor in $A/f_2A$, show that $f_1A \cap f_2A = f_1f_2A$.
We omit the straightforward argument.
\end{proof}

\noindent
Recall that we have defined the inverse image of a closed subscheme
under any morphism of schemes in
Schemes, Definition \ref{schemes-definition-inverse-image-closed-subscheme}.

\begin{lemma}
\label{lemma-pullback-locally-principal}
Let $f : S' \to S$ be a morphism of schemes. Let $Z \subset S$
be a locally principal closed subscheme. Then the inverse image
$f^{-1}(Z)$ is a locally principal closed subscheme of $S'$.
\end{lemma}

\begin{proof}
Omitted.
\end{proof}

\begin{definition}
\label{definition-pullback-effective-Cartier-divisor}
Let $f : S' \to S$ be a morphism of schemes. Let $D \subset S$
be an effective Cartier divisor. We say the {\it pullback of
$D$ by $f$ is defined} if the closed subscheme $f^{-1}(D) \subset S'$
is an effective Cartier divisor. In this case we denote it either
$f^*D$ or $f^{-1}(D)$ and we call it the
{\it pullback of the effective Cartier divisor}.
\end{definition}

\noindent
The condition that $f^{-1}(D)$ is an effective Cartier divisor
is often satisfied in practice. Here is an example lemma.

\begin{lemma}
\label{lemma-pullback-effective-Cartier-defined}
Let $f : X \to Y$ be a morphism of schemes.
Let $D \subset Y$ be an effective Cartier divisor.
The pullback of $D$ by $f$ is defined in each of the following cases:
\begin{enumerate}
\item $f(x) \not \in D$ for any weakly associated point $x$ of $X$,
\item $X$, $Y$ integral and $f$ dominant,
\item $X$ reduced and $f(\xi) \not \in D$ for any generic point $\xi$ of any
irreducible component of $X$,
\item $X$ is locally Noetherian and $f(x) \not \in D$ for any associated point
$x$ of $X$,
\item $X$ is locally Noetherian, has no embedded points, and
$f(\xi) \not \in D$ for any generic point $\xi$ of an irreducible component of
$X$,
\item $f$ is flat, and
\item add more here as needed.
\end{enumerate}
\end{lemma}

\begin{proof}
The question is local on $X$, and hence we reduce to the case
where $X = \Spec(A)$, $Y = \Spec(R)$, $f$ is
given by $\varphi : R \to A$ and
$D = \Spec(R/(t))$ where $t \in R$ is a nonzerodivisor.
The goal in each case is to show that $\varphi(t) \in A$
is a nonzerodivisor.

\medskip\noindent
In case (1) this follows from
Algebra, Lemma \ref{algebra-lemma-weakly-ass-zero-divisors}.
Case (4) is a special case of (1) by Lemma \ref{lemma-ass-weakly-ass}.
Case (5) follows from (4) and the definitions.
Case (3) is a special case of (1) by
Lemma \ref{lemma-weakass-reduced}.
Case (2) is a special case of (3).
If $R \to A$ is flat, then $t : R \to R$ being injective
shows that $t : A \to A$ is injective. This proves (6).
\end{proof}

\begin{lemma}
\label{lemma-pullback-effective-Cartier-divisors-additive}
Let $f : S' \to S$ be a morphism of schemes.
Let $D_1$, $D_2$ be effective Cartier divisors on $S$.
If the pullbacks of $D_1$ and $D_2$ are defined then the
pullback of $D = D_1 + D_2$ is defined and
$f^*D = f^*D_1 + f^*D_2$.
\end{lemma}

\begin{proof}
Omitted.
\end{proof}





